 \subsubsection{Backend}
 
\begin{figure}[H]
    \centering
    \includegraphics[width=1\textwidth]{SpecificaTecnica/Chat/VecDbBackend}
    \caption{Diagramma delle classi per il modulo di ricerca vettoriale: Backend}
\end{figure}

%TODO: decommentare questo
%Per la definizione di \textit{EmbeddedCartService}, vedere \ref{embeddedcartservice}.
 
\subsubsubsection{EmbedderPort}
 
\begin{figure}[H]
    \centering
    
\includegraphics[width=0.5\textwidth]{SpecificaTecnica/Chat/EmbedderPort}
    \caption{Ricerca Vettoriale: EmbedderPort}
\end{figure}
 
\paragraph{Descrizione}
Classe astratta Python che definisce la porta per la generazione di embedding testuali.\\
Segue il pattern architetturale delle porte e degli adattatori, disaccoppiando la logica
applicativa di \textit{EmbeddedCartService} ed \textit{EmbeddedCatalogService} dalla
specifica libreria di embedding utilizzata.\\
Tutte le classi adapter concrete devono implementare il metodo qui dichiarato.
 
\paragraph{Attributi}
Non definisce attributi propri, in quanto funge esclusivamente da interfaccia.
 
\paragraph{Metodi e funzioni}
\begin{itemize}
    \item \textbf{embed(text: string): np.ndarray} - restituisce il vettore di embedding
          corrispondente al testo passato in ingresso.
\end{itemize}
 
\subsubsubsection{VecDbPortOut}
 
\begin{figure}[H]
    \centering
    
\includegraphics[width=0.6\textwidth]{SpecificaTecnica/Chat/VecDbPortOut}
    \caption{Ricerca Vettoriale: VecDbPortOut}
\end{figure}
 
\paragraph{Descrizione}
Classe astratta Python che definisce la porta di output verso il database vettoriale.\\
Disaccoppia i servizi applicativi (\textit{EmbeddedCartService},
\textit{EmbeddedCatalogService}) dalla specifica implementazione del database vettoriale.\\
Le classi adapter concrete (\textit{CartVecDbAdapter}, \textit{CatalogVecDbAdapter})
devono implementare tutti i metodi qui dichiarati.
 
\paragraph{Attributi}
Non definisce attributi propri, in quanto funge esclusivamente da interfaccia.
 
\paragraph{Metodi e funzioni}
\begin{itemize}
    \item \textbf{reset()} - svuota completamente il database vettoriale, eliminando tutti
          i vettori e le associazioni precedentemente inserite.
    \item \textbf{search(v: np.ndarray, n: int, threshold: float): list[string]} -
          ricerca nel database vettoriale i prodotti il cui vettore è simile a \textit{v}
          entro la soglia indicata, restituendo al massimo \textit{n} identificativi.
    \item \textbf{add(prod\_id: string, vector: np.ndarray)} - aggiunge al database
          vettoriale il prodotto identificato da \textit{prod\_id} con il relativo vettore.
\end{itemize}
 
\subsubsubsection{SentenceTransformerEmbedder}
 
\begin{figure}[H]
    \centering
    \includegraphics[width=0.6\textwidth]{SpecificaTecnica/Chat/SentenceTransformerEmbedder}
    \caption{Ricerca Vettoriale: SentenceTransformerEmbedder}
\end{figure}
 
\paragraph{Descrizione}
Classe Python che incapsula il modello \textit{all-MiniLM-L6-v2} di
\textit{sentence-transformers} per la generazione di embedding testuali normalizzati.\\
Viene istanziata una sola volta e condivisa tramite \textit{EmbedderAdapter}.\\
Non implementa direttamente \textit{EmbedderPort}: tale responsabilità è delegata ad
\textit{EmbedderAdapter}.

\paragraph{Attributi}
\begin{itemize}
    \item \textbf{embedder: SentenceTransformer} - istanza del modello
          \textit{all-MiniLM-L6-v2} caricata al momento della costruzione.
\end{itemize}
 
\paragraph{Metodi e funzioni}
\begin{itemize}
    \item \textbf{embed(text: string): NDArray[np.float32]} - genera l'embedding
          normalizzato del testo, convertito in un array NumPy di tipo \textit{float32}.
\end{itemize}
 
\subsubsubsection{EmbedderAdapter}
 
\begin{figure}[H]
    \centering
    \includegraphics[width=0.6\textwidth]{SpecificaTecnica/Chat/EmbedderAdapter}
    \caption{Ricerca Vettoriale: EmbedderAdapter}
\end{figure}
 
\paragraph{Descrizione}
Classe Python che implementa \textit{EmbedderPort} seguendo il pattern degli adattatori.\\
Riceve un'istanza di \textit{SentenceTransformerEmbedder} alla costruzione e delega
interamente ad essa la generazione degli embedding.\\
Permette di sostituire la libreria di embedding sottostante senza modificare i servizi
applicativi.
 
\paragraph{Attributi}
\begin{itemize}
    \item \textbf{sentence\_transformer\_embedder: SentenceTransformerEmbedder} -
          istanza del componente concreto che esegue la generazione dell'embedding.
\end{itemize}
 
\paragraph{Metodi e funzioni}
\begin{itemize}
    \item \textbf{\_\_init\_\_(sentence\_transformer\_embedder: SentenceTransformerEmbedder)}
          - inizializza l'adattatore ricevendo l'istanza di
          \textit{SentenceTransformerEmbedder}.
    \item \textbf{embed(text: string): NDArray[np.float32]} - delega la generazione
          dell'embedding a \textit{sentence\_transformer\_embedder} e ne restituisce il
          risultato.
\end{itemize}
 
\subsubsubsection{EmbeddedCatalogService}
 
\begin{figure}[H]
    \centering
    \includegraphics[width=0.7\textwidth]{SpecificaTecnica/Chat/EmbeddedCatalogService}
    \caption{Ricerca Vettoriale: EmbeddedCatalogService}
\end{figure}
 
\paragraph{Descrizione}
Classe Python che coordina la ricerca vettoriale sull'intero catalogo prodotti.\\
Espone i metodi per caricare il catalogo nel database vettoriale e per cercare prodotti
semanticamente simili a una query testuale.
 
\paragraph{Attributi}
\begin{itemize}
    \item \textbf{CATALOG\_SEARCH\_LIMIT: int} - costante di classe che definisce il
          numero massimo di risultati restituiti dalla ricerca nel catalogo, impostata a
          30.
    \item \textbf{catalog\_vect: VecDbPortOut} - istanza del porto di output verso il
          database vettoriale del catalogo.
    \item \textbf{catalog\_repo: CatalogRepoPort} - istanza del porto di accesso ai
          dati del catalogo.
    \item \textbf{embedder: EmbedderPort} - istanza del porto per la generazione degli
          embedding.
\end{itemize}
 
\paragraph{Metodi e funzioni}
\begin{itemize}
    \item \textbf{\_\_init\_\_(catalog\_vect: VecDbPortOut, catalog\_repo: CatalogRepoPort,
          embedder: EmbedderPort)} - inizializza il servizio ricevendo le dipendenze
          tramite iniezione.
    \item \textbf{load\_catalog()} - recupera l'intero catalogo tramite
          \textit{catalog\_repo}, genera l'embedding di ciascun prodotto e lo aggiunge al
          database vettoriale.
    \item \textbf{search\_catalog(query: string, threshold: float): list[string]} -
          genera l'embedding della query e restituisce gli identificativi dei prodotti del
          catalogo semanticamente simili entro la soglia indicata.
\end{itemize}
 
\subsubsubsection{CartVecDbAdapter}
 
\begin{figure}[H]
    \centering
    \includegraphics[width=0.6\textwidth]{SpecificaTecnica/Chat/CartVecDbAdapter}
    \caption{Ricerca Vettoriale: CartVecDbAdapter}
\end{figure}
 
\paragraph{Descrizione}
Classe Python che implementa \textit{VecDbPortOut} seguendo il pattern degli adattatori,
specializzandosi per la gestione del database vettoriale del carrello.\\
Riceve alla costruzione un'istanza di \textit{FaissCartDb} e il nome utente corrente, e
trasparisce le operazioni sul database delegando a \textit{FaissCartDb} includendo il
contesto utente.
\paragraph{Attributi}
\begin{itemize}
    \item \textbf{faiss\_db: FaissCartDb} - istanza del database vettoriale FAISS
          dedicato al carrello.
    \item \textbf{username: string} - nome utente associato alla sessione corrente,
          utilizzato per filtrare i risultati della ricerca.
\end{itemize}
 
\paragraph{Metodi e funzioni}
\begin{itemize}
    \item \textbf{\_\_init\_\_(faiss\_db: FaissCartDb, username: string)} -
          inizializza l'adattatore ricevendo il database FAISS e il nome utente.
\end{itemize}
 
\subsubsubsection{CatalogVecDbAdapter}
 
\begin{figure}[H]
    \centering
    \includegraphics[width=0.6\textwidth]{SpecificaTecnica/Chat/CatalogVecDbAdapter}
    \caption{Ricerca Vettoriale: CatalogVecDbAdapter}
\end{figure}
 
\paragraph{Descrizione}
Classe Python che implementa \textit{VecDbPortOut} seguendo il pattern degli adattatori,
specializzandosi per la gestione del database vettoriale del catalogo.\\
Riceve alla costruzione un'istanza di \textit{FaissCatalogDb} e delega ad essa tutte le
operazioni senza aggiungere il contesto utente, in quanto il catalogo è condiviso.
 
\paragraph{Attributi}
\begin{itemize}
    \item \textbf{faiss\_db: FaissCatalogDb} - istanza del database vettoriale FAISS
          dedicato al catalogo.
\end{itemize}
 
\paragraph{Metodi e funzioni}
\begin{itemize}
    \item \textbf{\_\_init\_\_(faiss\_db: FaissCatalogDb)} - inizializza l'adattatore
          ricevendo il database FAISS del catalogo.
\end{itemize}
 
\subsubsubsection{FaissCartDb}
 
\begin{figure}[H]
    \centering
    \includegraphics[width=0.7\textwidth]{SpecificaTecnica/Chat/FaissCartDb}
    \caption{Ricerca Vettoriale: FaissCartDb}
\end{figure}
 
\paragraph{Descrizione}
Classe Python che gestisce un indice FAISS per la ricerca vettoriale nel carrello.\\
Mantiene una mappa \textit{user\_map} che associa ogni username alle posizioni dei propri
prodotti nell'indice, in modo da filtrare i risultati in fase di ricerca.
 
\paragraph{Attributi}
\begin{itemize}
    \item \textbf{dimension: int} - dimensione dei vettori dell'indice, impostata a 384
          per default.
    \item \textbf{index: faiss.IndexFlatL2} - indice FAISS basato sulla distanza
          euclidea L2.
    \item \textbf{prod\_ids: list[string]} - lista degli identificativi di prodotto
          ordinata per posizione nell'indice.
    \item \textbf{user\_map: dict[string, list[int]]} - dizionario che mappa ogni
          username alla lista delle posizioni nell'indice occupate dai suoi prodotti.
\end{itemize}
 
\paragraph{Metodi e funzioni}
\begin{itemize}
    \item \textbf{search(vector: np.ndarray, n: int, threshold: float, username: string):
          list[string]} - ricerca nell'indice i vettori più vicini a \textit{vector},
          filtrando i risultati per le sole posizioni associate a \textit{username} e
          restituendo gli identificativi con distanza L2 inferiore o uguale a
          \textit{threshold}.
\end{itemize}
 
\subsubsubsection{FaissCatalogDb}
 
\begin{figure}[H]
    \centering
    \includegraphics[width=0.7\textwidth]{SpecificaTecnica/Chat/FaissCatalogDb}
    \caption{Ricerca Vettoriale: FaissCatalogDb}
\end{figure}
 
\paragraph{Descrizione}
Classe Python che gestisce un indice FAISS per la ricerca vettoriale sul catalogo
prodotti.\\
A differenza di \textit{FaissCartDb} non mantiene una mappa per utente, in quanto il
catalogo è condiviso tra tutti gli utenti del sistema.
 
\paragraph{Attributi}
\begin{itemize}
    \item \textbf{dimension: int} - dimensione dei vettori dell'indice, impostata a 384
          per default.
    \item \textbf{index: faiss.IndexFlatL2} - indice FAISS basato sulla distanza
          euclidea L2.
    \item \textbf{prod\_ids: list[string]} - lista degli identificativi di prodotto
          ordinata per posizione nell'indice.
\end{itemize}
 
\paragraph{Metodi e funzioni}
\begin{itemize}
    \item \textbf{search(vector: np.ndarray, n: int, threshold: float): list[string]} -
          ricerca nell'indice i vettori più vicini a \textit{vector} e restituisce gli
          identificativi con distanza L2 inferiore o uguale a \textit{threshold}.
\end{itemize}

 \subsubsection{Backend}
 
\begin{figure}[H]
    \centering
    \includegraphics[width=1\textwidth]{SpecificaTecnica/Chat/CatalogBackend}
    \caption{Diagramma delle classi per il catalogo: Backend}
\end{figure}

\subsubsubsection{CatalogRepoPort}
 
\begin{figure}[H]
    \centering
    
\includegraphics[width=0.6\textwidth]{SpecificaTecnica/Chat/CatalogRepoPort}
    \caption{Ricerca Vettoriale: CatalogRepoPort}
\end{figure}
 
\paragraph{Descrizione}
Classe astratta Python che definisce la porta di accesso ai dati del catalogo prodotti.\\
Disaccoppia i servizi applicativi (\textit{EmbeddedCatalogService},
\textit{ToolCatalogService}) dalla specifica implementazione del repository.\\
Tutte le classi adapter concrete devono implementare i metodi qui dichiarati.
 
\paragraph{Attributi}
Non definisce attributi propri, in quanto funge esclusivamente da interfaccia.
 
\paragraph{Metodi e funzioni}
\begin{itemize}
    \item \textbf{get\_product(prod\_id: string): CatalogProduct | None} -
          restituisce il prodotto corrispondente all'identificativo indicato, oppure
          \textit{None} se non trovato.
    \item \textbf{get\_full\_catalog(): list[CatalogProduct]} - restituisce la lista
          completa dei prodotti presenti nel catalogo.
\end{itemize}

\subsubsubsection{CatalogRepoAdapter}
 
\begin{figure}[H]
    \centering
    \includegraphics[width=0.6\textwidth]{SpecificaTecnica/Chat/CatalogRepoAdapter}
    \caption{Ricerca Vettoriale: CatalogRepoAdapter}
\end{figure}
 
\paragraph{Descrizione}
Classe Python che implementa \textit{CatalogRepoPort} seguendo il pattern degli
adattatori.\\
Riceve alla costruzione un'istanza di \textit{CatalogRepository} e si occupa di
convertire i modelli ORM \textit{Anaart} negli schemi applicativi \textit{CatalogProduct},
mappando in particolare l'unità di misura tramite \textit{MeasureUnitEnum}.
 
\paragraph{Attributi}
\begin{itemize}
    \item \textbf{repo: CatalogRepository} - istanza del repository che accede
          direttamente al database.
\end{itemize}
 
\paragraph{Metodi e funzioni}
\begin{itemize}
    \item \textbf{\_\_init\_\_(repo: CatalogRepository)} - inizializza l'adattatore
          ricevendo il repository.
    \item \textbf{\_map\_measure\_unit(row: Anaart): MeasureUnitEnum} - metodo privato
          che converte il codice e la descrizione dell'unità di misura presenti nel record
          ORM nel corrispondente valore dell'enumerazione \textit{MeasureUnitEnum}.
    \item \textbf{\_map\_product(row: Anaart): CatalogProduct} - metodo privato che
          converte un record ORM \textit{Anaart} nello schema applicativo
          \textit{CatalogProduct}, delegando la mappatura dell'unità di misura a
          \textit{\_map\_measure\_unit}.
    \item \textbf{get\_product(prod\_id: string): CatalogProduct | None} - recupera il
          prodotto tramite \textit{repo} e, se presente, lo converte tramite
          \textit{\_map\_product}.
    \item \textbf{get\_full\_catalog(): list[CatalogProduct]} - recupera l'intero
          catalogo tramite \textit{repo} e converte ciascun record tramite
          \textit{\_map\_product}.
\end{itemize}
 
\subsubsubsection{CatalogRepository}
 
\begin{figure}[H]
    \centering
    \includegraphics[width=0.6\textwidth]{SpecificaTecnica/Chat/CatalogRepository}
    \caption{Ricerca Vettoriale: CatalogRepository}
\end{figure}
 
\paragraph{Descrizione}
Classe Python che gestisce l'accesso diretto al database per il catalogo prodotti.\\
Riceve una sessione SQLModel alla costruzione e la utilizza per eseguire le query tramite
il modello ORM \textit{Anaart}.\\
Rappresenta il livello più basso della catena
\textit{EmbeddedCatalogService} → \textit{CatalogRepoAdapter} → \textit{CatalogRepository}.
 
\paragraph{Attributi}
\begin{itemize}
    \item \textbf{db: Session} - sessione SQLModel attiva utilizzata per l'esecuzione
          delle query.
\end{itemize}
 
\paragraph{Metodi e funzioni}
\begin{itemize}
    \item \textbf{\_\_init\_\_(db: Session)} - inizializza il repository ricevendo la
          sessione del database.
    \item \textbf{get\_product(prod\_id: string): Anaart | None} - recupera il prodotto
          corrispondente all'identificativo indicato, restituendo \textit{None} se non
          trovato.
    \item \textbf{get\_full\_catalog(): list[Anaart]} - recupera tutti i prodotti
          presenti nella tabella \textit{Anaart}.
\end{itemize}

\subsubsubsection{ToolCatalogService}
 
\begin{figure}[H]
    \centering
    \includegraphics[width=0.7\textwidth]{SpecificaTecnica/Chat/ToolCatalogService}
    \caption{Ricerca Vettoriale: ToolCatalogService}
\end{figure}
 
\paragraph{Descrizione}
Classe Python che orchestra la ricerca semantica sul catalogo integrando la ricerca
vettoriale con le preferenze di acquisto dell'utente.\\
Recupera i prodotti tramite \textit{EmbeddedCatalogService} e li riordina privilegiando
quelli acquistati con maggiore frequenza dall'utente, qualora siano disponibili le
preferenze.
 
\paragraph{Attributi}
\begin{itemize}
    \item \textbf{embedded\_catalog: EmbeddedCatalogService} - istanza del servizio di
          ricerca vettoriale sul catalogo.
    \item \textbf{catalog\_repo: CatalogRepoPort} - istanza del porto di accesso ai dati del catalogo, utilizzata per recuperare i dettagli dei prodotti trovati.
    \item \textbf{preferred\_product\_frequency: dict[string, int]} - dizionario che mappa gli identificativi di prodotto alla frequenza di acquisto dell'utente corrente; valore di default è un dizionario vuoto.
\end{itemize}
 
\paragraph{Metodi e funzioni}
\begin{itemize}
    \item \textbf{\_\_init\_\_(embedded\_catalog: EmbeddedCatalogService, catalog\_repo:
          CatalogRepoPort, preferred\_product\_frequency: dict[string, int] | None)} -
          inizializza il servizio ricevendo le dipendenze; se
          \textit{preferred\_product\_frequency} è \textit{None}, viene sostituito con un dizionario vuoto.
    \item \textbf{search\_catalog(query: string, threshold: float): list[CatalogProduct]} - ricerca i prodotti semanticamente simili alla query tramite
          \textit{embedded\_catalog}, ne recupera i dettagli tramite
          \textit{catalog\_repo} e, se sono disponibili le preferenze utente, riordina i risultati privilegiando i prodotti con frequenza di acquisto maggiore, mantenendo l'ordine di similarità come criterio secondario.
\end{itemize}